可迁移的段落骨架目前归纳为一百五十类。第一类是观察段：数据或图形中出现了什么现象，该现象让
哪个假设或方法变得合理。第二类是改进段：原模型遗漏什么结构，这一遗漏会导致何种
偏差，因而补入哪个变量、交互项或约束。第三类是比较段：固定数据、指标和预算，
列出候选结果，再说明采用其中一个的依据。第四类是实验设计段：先写预算和安全
边界，再列候选方法、舍弃理由、因素水平与最终组合。第五类是成对结果段：在同一
对象下先给基准或无约束结果，再给受限或调整结果，最后解释差值对决策的含义。
第六类是小样本裁决段：先报看似占优的指标，再交代样本数、自由度或外推限制，最后
说明为何保留较低阶方案。第七类是递进指标段：复述上一轮指标和剩余不足，写出本次整理
新增结构，再用新指标决定继续、停止或换用另一类方法。第八类是双响应控制变量段：
先写保持条件、唯一改变量和比较组，再分别读两项响应，只作局部判断。第九类是诊断
修模段：初始模型先点名多重共线性等具体问题，给出逐步回归等修正动作，随后按相对
拟合、绝对误差和整体检验三个口径解释修正结果。第十类是潜变量代理段：先说明业务概念
不可直接观测，再列代理指标、方向和权重解释。第十一类是替代量否决段：先写方便的候选
口径，再用两个不重叠的理由说明它会丢失什么信息，最后给替代构造。第十二类是反常结果
修正段：结果与预期不符时先诊断机制，只修改受影响的约束或系数。第十三类是状态传递段：
明确本期输入、决策、实现量和期末状态如何形成下一期初值。第十四类是宏观--微观解释段：
总体序列回答方案是否闭合，代表时点说明状态和约束如何具体作用。第十五类是有限效果段：
同时给最大值与平均值，直接承认改善幅度有限，并解释可行域或资源边界。第十六类是模型沿用
段：逐项核对输入、输出和目标是否改变，再决定复用、修改、删除或新增模块。
第十七类是规则筛样段：先说明每条阈值分别排除何种不稳定，再规定同时满足后的统计量。
第十八类是失败驱动改权重段：点名固定权重造成的局部缺货状态，再由当前缺口更新下一轮权重。
第十九类是数据定窗段：用连续淡季或峰谷位置确定前视/后视长度，并解释库存调配方向。
第二十类是同口径裁决段：当成本、收益或误差对应的产出规模不同时，先改成单位产出指标再比较。
第二十一类是非对称损失段：先写两类偏差的业务后果差异，再列函数必须满足的方向与曲率条件，
最后构造分段表达式。第二十二类是观察后舍弃段：先说明规律和可能成因，再用干扰来源与任务变化
解释为何不纳模。第二十三类是初值与可分性段：先判断系统是否已有历史状态、过程是否有时间记忆，
再决定初值来源和能否逐期拆分。第二十四类是保守聚合段：先明确局部解的业务语义，再说明跨场景
应取最大还是最小，以及该方向怎样保证统一可行。第二十五类是反证假设段：提出便利假设，
再代入容量、数量或需求事实得到矛盾，最后修正假设。第二十六类是权重异常修正段：先给客观权重，
点名与业务不符的指标，再限定主观修正幅度并做一致性检查。第二十七类是贪心停止段：先定义满意
边界，达到即停，并说明继续选择会增加的采购、运输或库存代价。第二十八类是候选机制裁决段：
所有候选共享历史基准，分别说明变动机制，再按约束、目标、初期稳健性和现实含义逐个淘汰。
第二十九类是随机代表段：报告重复运行，只以频率和目标双条件选典型方案，不宣称结果唯一。
第三十类是目标逆转扩围段：先比较新旧目标，再解释先前低排名对象为何重新有贡献，据此调整候选范围。
第三十一类是分类压缩段：先列完整情形，随后指出可旋转、平移或交换重合的等价类，最后给保留分支和
适用条件。第三十二类是计算动作段：按固定对象、遍历对象、范围/步长、记录项和图表产物交代仿真，
使代码循环在正文中有清楚语义。第三十三类是长结果分流段：说明完整结果的存放位置、正文选取的代表
情形和代表理由，避免以“篇幅有限”掩盖证据。第三十四类是结构转化复用段：先完成坐标、图结构或
子问题分组的转化，逐项核对旧模型接口，随后指出哪些推导不重述、哪些特殊对象仍需另建关系。
第三十五类是负担触发换模段：先保留旧方法及其分类或计算负担，随后指出仍可利用的不变量，最后
给出新构造和被消去的工作量。第三十六类是数值到区间证明段：由具体例子定位可能区域，再检查
轨迹端点、给出变量界，并用互不交叠的区间或表格闭合穷举。第三十七类是局部让步--全局反转段：
先承认某算法在一轮或一个局部指标上占优，再说明状态更新或优先处理大误差点怎样改变总体目标。
第三十八类是任务充分停止段：同时列数值上还能达到的精度、任务阈值、额外轮次或资源代价，说明
为何停止在满足任务而非机器精度的位置。第三十九类是复用前失效审计段：先说明新问本质上复用
旧模型，再集中列出共线、共圆或接口退化，逐项给检测条件和修复动作。第四十类是信息边界账本段：
先写对象之间不共享什么，再列局部可观测量和局部目标；若需要全局指标，只把它作为事后评价，并
解释为何没有进入对象决策。第四十一类是预处理双收益段：先证明误差界缩小，再说明搜索区域、函数
评估或信号次数怎样减少。第四十二类是双轴或三轴方案裁决段：在同一初始数据下依次比较收敛速度、
稳定误差和资源代理，公开代理口径后再选择。第四十三类是引理到流程装配段：逐条写固定点、调整点、
控制角/比例和几何目标，再映射到具体对象与阶段。第四十四类是单样本不足段：先承认一组初值不足以
检验稳定性，再定义扰动分布、重复数据集、收敛区间和结论范围。第四十五类是缺失语义裁决段：
先判断空值在题面中表示未检出、未知还是未采样，再用同条件多值性和采样范围决定置零、拒绝填补或
局部忽略。第四十六类是检验前提分流段：先列候选检验的期望计数等前提统计量，再按满足与否选择
Pearson、Yates 或其他方法。第四十七类是双向证据段：用前问趋势和已知标签解释后问聚类，再用
聚类结果反向核对前问结论。第四十八类是指标门控计算段：逐对象报告 $R^2$ 等指标，把高低分支
分别接到预测、保持不变或改用基线。第四十九类是异常距离局部修正段：先点名唯一敏感量和异常放大
造成的失真，再只对该量作单调压缩，不借此替换整条主模型。第五十类是变量--样本双聚类段：R 型
聚类先降相关并选代表变量，Q 型聚类后分样本，分别说明两个对象和输出。第五十一类是数据语义删步
段：比较新旧表单的采样或标签语义，核对新表已直接给出状态后，删除重复估计而复用后续模型。
第五十二类是同基准类型比较段：固定一个参考序列，分别报告各类最大、最小关联及集中/离散程度，
最后解释类型差异。第五十三类是同字段分组填补段：先列匹配条件和候选样品，再按参考质量选择
热卡或众数，并写出具体编号与填充值。第五十四类是近似零值辨析段：区分空白和报表零的来源，
说明为何不能作普通零，再接乘法替换与闭合。第五十五类是主导成分遮蔽拆图段：保留全部图，
另画去除主导成分的图，结合两图读取大小量变化。第五十六类是可解释规则--多变量复核段：
简单阈值先给直观分类和测试结果，PLS-DA/VIP 再检查多变量差异与一致性。第五十七类是无监督
身份判定段：没有已知标签或因变量时，先说明不能监督分类，再用曲线选择聚类数。第五十八类是
聚类强弱诚实限定段：并列两类分簇图，一类清楚便具体描述，另一类不清楚只作粗分并限制结论。
第五十九类是双指标定维段：同时比较交叉验证误差与累计解释率，选择最小充分维数。第六十类是
对称矩阵配对检验段：先由对称性取上三角去重，再由成分对对应关系选择配对检验并报告决策。
第六十一类是分问关系定性段：先说明研究对象与本质身份，再依次说明正向计算、反问题和放宽约束。
第六十二类是关键困难隔离段：从题面公式中剔除可直接获得的量，只把剩余未知量交给后续计算模块。
第六十三类是对称与统计补偿简化段：先写精确边界的负担，再由数量占优、中心对称或统计补偿推出
中心判据和二值动作。第六十四类是坐标转移判定段：对象由 A 坐标转至公共坐标，再转入 B 坐标，
通过平面交点和区域内外完成有效/无效判断。第六十五类是能量收支闭合段：从全反射能量扣除无效点
损失，再以接收能量或射线数形成效率比值。第六十六类是空间证据到布场段：由云图或方向性读出
高效区域，把镜面集中方向、塔址反向移动、特征距离和分区层规则连续写出。第六十七类是计算预算到
粗细搜索段：先说明直接高精度遍历为何过重，再由二分或大步长定位初步范围，缩短步长精化并给停止
精度。第六十八类是范围--方向--饱和--排序段：明确扰动对象与步长，报告变化区间、增减非对称、
饱和位置和分指标敏感度次序。第六十九类是过程分段摘要：把物理或业务链拆成连续阶段，每一阶段
使用一个动作动词并交付一个下游量。第七十类是等价几何准备段：先给分解关系，随后将原判定改写为
可计算的子判定，最后说明二者为何等价。第七十一类是近邻栅格仿真段：先限定候选近邻，再写射线、
平面和区域判定，避免把全场两两检查塞入一句。第七十二类是快速核误差门段：以高精度计算为基准，
说明低分辨率核节省什么，并用同一输入报告偏差后限定其用途。第七十三类是固定其余变量降维段：
列出固定量，只让一个变量改变；从响应曲线确认形态后，再选择一维搜索。第七十四类是空间图触发
方向变量段：由散点或云图读出高效方向，把方向转成设施移动、坐标符号和搜索区间。第七十五类是
继承旧解并开放新自由度段：先说明前问已经确定的变量，再逐项冻结，只把新增自由度交给本问求解。
第七十六类是结构预言检验段：先由机理写出图形或峰值应怎样变化，主动改变输入结构，重新计算后
以观察到的结构与预言是否一致作裁决。第七十七类是条件分支传播段：先写第一阻断条件，成立时
明确跳过的后续计算；不成立再进入第二、第三判定，只有全部通过才计算接收或输出。第七十八类是
邻接降算量段：先给全对全对象数、网格与角度共同形成的规模，再由物理局部性设距离阈值并只遍历
候选集合。第七十九类是抽样粗细寻优段：小比例多次抽样负责估计粗参数，冻结弱变量后提高比例、
减小步长并承担最终评价。第八十类是弱影响变量冻结段：固定基线分别试算变量，报告指标变化区间，
取代表或较优值后把搜索预算转向关键变量。第八十一类是趋稳阈值选步长段：列出分辨率序列和指标，
指出增幅开始变缓的位置，用“接近稳定”限定结论，再给工作步长。第八十二类是达标结果反推目标段：
方案超过硬阈值后，分析删去低效或高干扰对象对总量、面积、干扰和单位量的连锁作用，重新决定
次目标并验算可行性。
第八十三类是唯一剩余量隔离段：把题面数据代入既有模型，逐项列出可由题面求得的量，只留下唯一不能由现有量
求得的量，解释随机或高维原因后再给抽样法与样本数。第八十四类是重叠区域离散段：先指出两区域
可能重叠、直接相加会重复计数，再给离散对象、网格步长和并集计数规则。第八十五类是空间猜想
复核段：由序号或时间曲线提出位置猜想，另选代表日期和时刻绘制二维分布，复核后才转成布局动作。
第八十六类是方向抵消降维段：固定其余参数沿两个方向分别扰动，依据对称性、跨时段抵消和影响
幅度冻结弱方向坐标，只保留主方向。第八十七类是前问结构按组降维段：前问得到规则结构，本问
先说明沿用，再由对称性让同组对象共享参数，把逐对象变量压成逐组变量并重新选择求解器。
第八十八类是可行样本背景检验段：在相同硬约束下生成 $N$ 组可行基线，以同指标、同单位比较，
只陈述选定方案在该样本背景中的优势，不外推为全局最优。
第八十九类是两个关键量归约段：先写旧模型的输入输出接口，点名两个缺失几何量，分别推导后
代回旧模型，不用“沿用前问”省略接口。第九十类是命名对象后推导段：先命名平面、直线和向量，
再写位置/垂直关系、叉乘、坐标表达和物理意义。第九十一类是规范加自我核验段：规范说明常用做法，
命题把最优性压成同长度候选比较，覆盖面积或同口径指标承担裁决。第九十二类是边界锚定贪心段：
选起始边界，以恰达边界确定首个对象，按固定空间步长保留可行候选，取最低可行重叠率并迭代至
对边。第九十三类是绝对--相对双误差检验段：基线与随机对照比较同一位置或指标，同时报告绝对
误差和相对误差，只作同计量方法一致性结论。第九十四类是复杂度触发换求解器段：先列未知量个数、
取值范围和遍历成本，再给群智能算法、参数和停止规则。第九十五类是残差驱动区域细化段：按结构
初分区域，拟合局部代理并读取 MSE，只拆分失败区、重拟合子区而保留合格区。第九十六类是分区域
结果解释段：点名异常区域和指标，用局部坡度、覆盖宽度或作业简化解释线数、间距和漏测后果。
第九十七类是竞争定义裁决段：承认题意歧义，保留两个候选定义，固定共同数值例，寻找可区分的
观察量，再写保留与舍弃。第九十八类是改变量复用段：比较前后两问，列不变结构和改变量，只修改
受影响的函数输入。第九十九类是小参数长尺度累积段：把局部小量换算到完整任务尺度，以累计差值
和两端可行区间判断能否忽略。第一百类是内部裁决加外部印证段：先由题内数值比较作选择，再说明
结果与标准或规范一致。第一百零一类是四项结构检查段：依次核对对称、单调、特殊方向不变和极限式，
每一项对应独立预测。第一百零二类是对边闭合递推段：从首边界和首个对象开始，按固定规则递推，
计算最后可行对象的对边余量，并舍弃越界候选。第一百零三类是读者连续性插图段：后文确需调用
早先几何时，说明为减少回翻而就近重现，只保留在当前推导中所需的标记。第一百零四类是最不利位置全局
覆盖段：定位等值线间距最宽或最平缓处，说明其最不利性，由该处充分条件推出全域覆盖并核边界。
第一百零五类是侧写到决策接口段：逐项写侧写对象、观察、所支持的变量/窗口/约束和后续动作，
不让探索性图表与模型脱节。第一百零六类是代表规则先行段：对象过多时先给排序指标、阈值和代表
选取规则，再展示代表结果并标出全部附件位置。第一百零七类是系数业务量纲段：从回归系数读出
自变量每变化一个业务单位时因变量的方向和变化量，再说明经营含义。第一百零八类是特征语义账本段：
依次写特征对象、计算口径、代理现象和下游用途，不以特征名称代替构造理由。第一百零九类是现有--
外部证据分流段：先列附件内能够量化的判断，再列尚未覆盖的信息、采集方式和预定用途。第一百一十类
是代理竞争机制段：同一代理量保留至少两个可行解释，指出能够区分它们的追加证据，并分别给出动作。
第一百一十一类是主文--附件规则分工段：正文保留选择规则、代表结果、核心数值和解释，附件承接
全部表、补充图及代码，且正文明确给出去向。第一百一十二类是数据病态触发换口径段：依次写
分布、零值、时间结构和计算负担，再给转换规则、新数据口径与下游模型。第一百一十三类是判断矩阵
压缩段：先定义阈值、显著性和字符串含义，再以 T/F 等接口代替正文逐格复述，并给全部系数去向。
第一百一十四类是业务时间尺度授权段：从保鲜、余货或运营时滞得到记忆长度，由此确定阶数或窗口，
不以“综合考虑”跳过尺度推导。第一百一十五类是检验表导读段：先说明各行列对应估计量、不确定性
和检验量，再按信息准则、拟合和显著性分别裁决，最后限定任务与时窗。第一百一十六类是异常数值
业务翻译段：把负值、零值或极端量先映射为库存、退货、折价等可观察状态，再决定保留、截断或转成
动作。第一百一十七类是低成本独立核对段：主模型给出候选后，用频率、阈值或业务规则独立筛选，
比较名单和差异，只在口径可比时称有限一致。第一百一十八类是预测--优化中间量接口段：逐项写
预测量进入目标或约束的位置，并把近期均值、上期状态、损耗、成本和业务下限接到下一阶段决策。
第一百一十九类是开放题聚焦段：先说明题面没有指定探究方向，随后将对象拆为能够互相补足的观察面，
并交代这些观察面共同覆盖什么任务。第一百二十类是候选失败换模段：列出实际试过的同类方法、当前
指标或现象上的不足，再说明新方法正好处理哪项困难。第一百二十一类是数据能力触发降级段：写清
缺失、短窗、无趋势或随机波动为何使复杂模型不可靠，再给简单估计及其时窗边界。第一百二十二类是
变量结构授权求解器段：比较前后问的连续、离散和多目标结构，只在变量域或目标关系改变后切换算法。
第一百二十三类是弹性修正接口段：先给基线预测，再用响应或弹性函数修正，逐项指出修正量、成本、
损耗和补货量进入目标或约束的位置。第一百二十四类是同期历史基线段：先解释为何选同季、同星期或
同业务状态，随后将方案数值放入该参照范围作有限比较。第一百二十五类是从模型边界生长建议段：先写
当前模型由哪些内部信息驱动，随后指出它不能区分的现实关系，最后列能区分关系的数据和动作。第一百
二十六类是扰动范围限定段：写明被扰动输入、范围、其余保持量和观测指标，只把结论限定为当前扰动
与报告范围内的稳定性检查。
第一百二十七类是局部递推展开段：先求起点，再定义从第 $i$ 个对象到第 $i+1$ 个对象的关系，最后说明如何沿空间与时间铺开。第一百二十八类是物理条件选根段：先列顺序、方向或单调条件过滤数学根，再以局部连续性选最近可行根。第一百二十九类是首次事件缩域段：由结构重复、轨迹继承或单调关系限定最早事件涉及的对象，并说明为何不会漏掉更早情形。第一百三十类是事件顺序消枝段：算法遇到首次事件即停止时，明确哪些只会随后发生的分支无需进入当前判定。第一百三十一类是越界回退求精段：写清初始步长、前进判据、越界回退、缩步比例、停止阈值和由此得到的误差界。第一百三十二类是固定量否决优化段：先由相切、守恒或代数约束证明目标量唯一，说明在当前约束下没有优化自由度，再进入指定方案。第一百三十三类是跨区域复用段：复用前问公式时逐项列出不等式、坐标符号、切向量、初值和边界的改变。第一百三十四类是齐次传递缩放段：先证明局部状态改变 $k$ 倍会传到相邻状态，再用已有解与上限求整体比例。第一百三十五类是同输出灵敏度比较段：对多个输入分别作同尺度扰动，始终观察同一输出，最后按变化方向和幅度给敏感度排序。
第一百三十六类是反证逐节前推段：先假设第 $i$ 个构件最先触发事件，再由轨迹滞后推到
$i-1$，逐节前推至首构件并否定假设。第一百三十七类是事件对象切换解释段：参数微变而
输出跃变时，先核对首次触发对象是否改变，再用局部结构差异解释不连续。第一百三十八类是
终点可行反例段：对只检查终态的简化方案，沿全过程搜索更早事件，以一个明确反例裁决其
适用性。第一百三十九类是现实动作推出边界段：说明参数继续变化会要求哪个构件执行怎样的
不可能动作，再据此给出上下界。第一百四十类是异常时段局部解释段：先圈定异常时段，再比较
同一时间内不同节点经过的路程，由局部轨迹解释速度或响应变化。

